In-field nutrient management remains the primary lever for reducing nitrate
and phosphorus (P) losses from tile-drained land, yet the levers themselves are
better described than optimised. The frontier is less about which practices
help than about the metrics, mechanisms, crediting rules, and temporal targets
that would let managers set the right rate, form, and timing for a given soil,
climate, and cropping system. \textcite{Ma2025} frame the programme broadly,
arguing that far more work is needed to understand the tradeoffs among
productivity, nutrient-loss reduction, and ecosystem services in tile-drained
systems. This section synthesises the open questions for nitrogen, for
phosphorus, and for cropping-system change; model-algorithm gaps,
greenhouse-gas tradeoffs, edge-of-field treatment, and spatial bundling are
treated only in passing here and developed elsewhere in the review.

\subsection{Refining nitrogen management and the timing of loss}

Refining fertiliser nitrogen (N) is the most consistently identified priority.
\textcite{ROSSI1991} call for careful evaluation of fertiliser timing, rate,
form, and placement, and \textcite{Zhao2000} for strategies that raise N-use
efficiency in corn through optimal rate, timing, and form. \textcite{Randall2001}
add crediting of previous legumes and manure and improved manure management,
while \textcite{Dattamudi2020} conclude that application schemes need
reevaluation toward lower corn rates, precision agriculture, and split, properly
timed applications---echoing \textcite{Elmi2002}, who see reductions from
matching rate and timing to crop uptake. Optimising N rates is expected to
retain primary research and regulatory focus, with accurate 4R implementation
requiring site-specific knowledge to balance agronomic and environmental goals
\parencite{Christianson2015b}.

\subsubsection{Next-generation rate tools and decision support}

\textcite{Bandura2025} propose concrete standards for the next generation of
N-rate tools: developed against the economically optimum N rate (EONR),
adequate if they predict need within 10~kg~N~ha$^{-1}$ of EONR, and
environmentally unacceptable if they exceed it by more than 30~kg~N~ha$^{-1}$ at
an N:corn price ratio of 5.6. \textcite{Liu2021} call for recommendations that
weigh environmental and economic benefits, case-by-case cost analyses, and
coordinated agency action to promote 4R adoption, and \textcite{Banger2018} for
experiments across wider soils and management to identify yield tradeoffs from
rate reductions---shifting emphasis toward cover crops and edge-of-field
practices where yield-loss risk appears. Comprehensive tools integrating weather
forecasts, management history, and local field parameters, with regional
agronomic-optimum-N-rate benchmarks \parencite{Maas2022}, precision application
by soil-N availability \parencite{Gupta2023}, better crediting of non-fertiliser
N and variable-rate models incorporating remote sensing \parencite{Dinnes2002},
and low-cost proxy variables for soil N-supply capacity \parencite{Cassman2025}
are all advanced. Local decision approaches also need refinement: late-spring
nitrate-test programmes are framed as a start to be replaced as research emerges
\parencite{Jaynes2004}, with grain-protein monitoring and risk insurance for
adopters \parencite{Karlen2005} and weather-based tools to reduce soil-testing
needs \parencite{Malone2010}. \textcite{Jaynes2013} seek to pinpoint where
between the V12 and V16 stages sidedress N ceases to benefit, and
\textcite{Jaynes2015} note sidedressing can be paired with soil tests or crop
N-deficiency sensing for real-time adjustment; comprehensive volatilisation
studies and data from data-poor regions are further gaps \parencite{Chatterjee2020b}.

\subsubsection{Timing, form, and application method}

Method-level questions remain open. \textcite{Malone2007b} highlight quantifying
split versus single spring applications, and \textcite{HernandezRamirez2011}
note it is unknown whether reduced-rate fall manure could match sidedress
environmentally, pointing to nitrification inhibitors and more complex
rotations; \textcite{Elmi2005} stress quantifying all major N loss pathways
simultaneously. Product-level frontiers include enhanced-efficiency fertiliser
with a nitrification inhibitor tested alongside drainage water management
\parencite{Kaur2024c}, NH$_4^+$-N stabilisers and cover-crop effects on
NH$_4^+$-N dynamics \parencite{Vincent2025}, and fertigation of silage corn with
its crop-injury and pathogen tradeoffs \parencite{Gamble2018}. Evening out
manure application in high-livestock catchments \parencite{Bechmann2014} and
comparing N-use consequences under global-warming projections
\parencite{Haghnazari2020} round out this cluster.

\subsubsection{Crediting residual, mineralisable, and legacy nitrogen}

A recurring, more specific gap is crediting N supplied by soil organic matter,
prior legumes, manure, and land-use legacies. \textcite{Bakhsh2007,Bakhsh2005}
call for assessing proper credit from mineralisation of the organic-matter pool
and soybean residue before applications, and \textcite{Randall2001} for improved
soil-N tests to estimate mineralisable and carryover N---echoed by a spring test
after dry years \parencite{RANDALL1995} and pre-sowing mineral-N measurement
\parencite{Conrad2009}. Systems transitioning from perennial cover need
second-year credits after grass conversion \parencite{Huggins2001} and care
after alfalfa \parencite{ROBBINS1980}, while irrigated vulnerable soils call for
reduced early-year N with residual-nitrate crediting \parencite{Derby2009}.
Legacy N is itself underexplored: a clearer definition is needed to apportion
losses \parencite{Williams2022}, and the coupled biogeochemical and hydrological
processes regulating legacy-N formation and its effect on nitrate export require
dedicated study \parencite{Hu2024}. Reconsidering fall application under warmer
winters \parencite{Gentry2014}, resin-membrane strips to characterise soil-N
cycling \parencite{Masri2024}, and better long-term flow-weighted-N predictions
for benchmark soils \parencite{Skaggs2012} complete the agenda.

\subsubsection{The temporal window of loss}

A distinctive frontier is the temporal concentration of losses and the need for
management targeted to that window. \textcite{Bakhsh2007} argue strategies must
reduce early-spring nitrate leaching, and \textcite{Bakhsh2010} extend this to
March--June while flagging the temporal distribution of growing-season rainfall.
Assessment of nitrate status and tile loss during April, May, and June warrants
study \parencite{Nance2007}, and a winter-hardy cover crop ahead of corn with a
shift of fall N to spring may best cut loss in intensively tiled watersheds
\parencite{Gentry2024}. For outlet management, \textcite{Skaggs1995} recommend
selectively managing outlets to reduce peak nitrate losses, especially after
drought-reduced yields, and \textcite{Elmi2004,Elmi2002} note that keeping
controlled drainage operational in early spring and late fall---when losses are
severe and drainage is not needed for production---could achieve further gains.

\subsection{Phosphorus metrics, mechanisms, and transport}

Phosphorus in subsurface drainage is repeatedly flagged as under-researched
relative to N. \textcite{Hertzberger2019} call for greater effort on P transport,
and \textcite{Yuan2022} for research on conservation-practice effectiveness in
reducing dissolved-P loss and the causes of its annual variability.
\textcite{Pease2018} identify pedologic controls on P transport as key needs,
and \textcite{deOliveira2022} recommend refining loss frameworks through
within-field studies and a better understanding of the ``where'' and ``when'' of
drainage P loss. A first-order requirement is to resolve partitioning and
sources: \textcite{King2015} call for enhanced P routines and simultaneous
measurement of P losses in surface runoff and subsurface drainage to determine
partitioning, \textcite{Williams2016} for differentiating soil versus fertiliser
sources of tile-drained P, and \textcite{Qi2018} for measuring the size
distribution of particulate-P carriers.

\subsubsection{Beyond agronomic soil-test P: saturation and sorption metrics}

A consistent theme is that agronomic soil-test P (STP) is insufficient for
environmental risk. \textcite{Sharpley1995} argue that environmental tests
estimating bioavailable P and sorption capacity should replace agronomic tests
on high-risk soils. \textcite{Hooda1999} propose that degree of P saturation
(DPS) may better estimate loss potential than extractable P such as Olsen-P, and
\textcite{Nguyen2002} call for critical DPS values calibrated across sediment-P
tests and equilibrium P concentrations. \textcite{Osterholz2020} recommend
assessing fields in the 100--400~mg~kg$^{-1}$ STP range and examining sorption
saturation, macropore formation, and mineralogy, while P-adsorption capacity
should precede developing organic soils for cropping \parencite{MILLER1979}.
\textcite{McDowell2001} propose that P isotope-exchange kinetics could explain
and eventually predict loss across a wide range of STP.

\subsubsection{Desorption, colloidal, and particulate transport}

The sorption/desorption interface and particulate pathways are poorly resolved.
\textcite{Mumbi2025} regard quantifying when and how varying water sources and
fluxes desorb and adsorb P as critical to predicting subsurface DRP transport,
and \textcite{Flores2021} note the need to understand particulate settling and
P desorption. \textcite{Uusitalo2003} argue that mitigating bioavailable
particulate P requires reducing particle mobilisation across the whole field
surface and that controlling diffuse DRP is a longer-term task than reducing
erosion, with the availability of eroded P to algae still limited
\parencite{Ulén2007}. \textcite{Wilson2019} highlight quantifying the colloidal-P
contribution to DRP, and \textcite{Norgaard2013} recommend including both
dissolved and colloid-associated leaching parameters in field-scale risk
assessment. Sediment-P release under alternating waterlogged--drying regimes
\parencite{Nguyen2002}, P fractionation under polyacrylamide \parencite{Jiang2010},
and combining columns, plots, and catchments to link particulate loss to
macroporosity \parencite{Laubel1999} are further priorities.

\subsubsection{Management effects, events, and legacy phosphorus}

\textcite{Uusitalo2018} note that assessing the pace of P stratification needs
more data and leave open whether no-till adoption as a eutrophication measure is
justified. \textcite{Ni2020} call for studies of optimal STP under changing
climate to address STP--precipitation interactions, and \textcite{Wang2018} ask
whether pre- or post-application tillage decreases tile DRP and whether findings
extend to manure injection, with immediate post-application manure risk still
poorly understood \parencite{Madison2014}. \textcite{Brendel2019} recommend
designing practices that reduce P losses during frequent small events, and
\textcite{Schilling2017b} refining organic-P/total-P ratios and quantifying
wastewater-treatment contributions. \textcite{Jarvie2017} ask how much of a 40\%
reduction target 4R practices alone could achieve and warn that drawing down
legacy P may be increasingly costly and intrusive; cost-effective hierarchies
prioritising lower Olsen P and targeted alum \parencite{McDowell2015} and P
removal by grass harvesting \parencite{Zhang2015b} are advanced. Fragmented
reporting also hampers synthesis: \textcite{Christianson2016} call for field
studies documenting multiple P species and further developing the Drain Load
database, \textcite{Macrae2019} for understanding hydroclimatic event drivers
and tillage or placement effects, \textcite{Tanga2025} for monitoring
particulate and dissolved forms with site-specific stoplog depths,
\textcite{Zhang2017} for new soil-P testing under no-till, and \textcite{Spiese2023}
for novel practices addressing nitrate and DRP simultaneously.

\subsubsection{Phosphorus in organic soils}

Phosphorus in cultivated organic soils is singled out as poorly understood.
\textcite{Grenon2022} call for mechanistic studies of the pathways, hot spots,
and hot moments of P transfer from agricultural organic soils, and
\textcite{Grenon2021} propose a concrete agenda using isotope tracing, XANES,
and $^{31}$P~NMR alongside sequential fractionation, microbial-turnover
assessment, and long-term drainage-water-management studies. Nutrient leaching
from abandoned peat fields, and the true nature of H$_2$SO$_4$-P, remain poorly
studied \parencite{Höyhtyä2025}.

\subsection{Cover crops, perennials, and rotational diversification}

\subsubsection{Cover crops: agronomic, environmental, and modelling frontiers}

Cover crops---especially winter rye---attract the largest concentration of open
questions. \textcite{Kladivko2014} call for research and technical assistance to
overcome adoption obstacles and for field experiments across more Midwestern
soils, and \textcite{Birru2023,Zhao2016} stress that existing evidence is drawn
from too few, geographically clustered sites and must be extended across more
species, soils, and climates. Agronomic priorities include tools to promote rye
in colder regions \parencite{Waring2020}, establishment strategies under
temperature and precipitation extremes \parencite{Bawa2023}, evaluation of more
cultivars to reduce yield variability \parencite{Kaspar2007}, breeding for cold
hardiness \parencite{Dinnes2002,Liu2019}, whether yield drag reflects poor
planter performance through residue \parencite{Dougherty2020}, and better
management of kura clover as living mulch \parencite{Qi2011b,Qi2011}. On
environmental mechanisms, authors seek research on the spring release of
scavenged N and its synchronisation with crop growth \parencite{Lacey2015},
adjustable seeding rates \parencite{Roth2018}, combined cover-crop and manure
effects on subsurface DRP \parencite{Askar2023}, and optimisation of cover-crop
services against disservices to the water balance \parencite{Meyer2019}.
\textcite{Liu2019} call for coordinated common experiments to quantify
winter-cover P contributions, longer-term studies remain scarce
\parencite{Kaspar2012,MartinezFeria2016}, growing cover crops longer in spring
may prevent N$_2$O increases \parencite{PrezaFontes2023}, and because over-winter
effects differ for N versus P, practices should target individual pathways
\parencite{McNeill2026}---with \textcite{Krueger2013} recommending practices
beyond rye where subsurface nitrate stayed high and rye had no effect on DRP.
Cover-crop modelling is a distinct frontier: validated rye growth and plant-N-
uptake algorithms \parencite{Feyereisen2006,Malone2007c}, dedicated winter-rye
modules and calibration \parencite{Malone2014b,Chatterjee2025}, and correction
of cover-crop and frozen-soil weaknesses \parencite{Dietzel2016,Gupta2022,Gupta2023b}
are sought, with detailed algorithm gaps developed in the modelling section.

\subsubsection{Perennials, bioenergy, and rotational diversification}

Perennial and land-use-change options feature prominently. Poorly drained
depressions are proposed for perennial bioenergy crops such as
\textit{Miscanthus $\times$ giganteus} or wetland restoration \parencite{Huang2023},
and keeping soil covered year-round with perennial ley is judged the most
effective way to mitigate N losses, with the P focus on decreasing surface
runoff \parencite{Norberg2022}. Additional data on alfalfa's effect on subsurface
discharge during large events \parencite{Arrueta2025}, a clear causal link
between perennial crops and nutrient loads with policies rewarding environmental
services \parencite{Arrueta2022}, robust water-quality datasets for perennialised
systems \parencite{Wilson2023}, and large-scale bioenergy-scenario assessment
with cost--benefit analysis to alleviate Gulf hypoxia
\parencite{Guo2018,Woodbury2018} extend the agenda. Rotation and diversification
questions recur: accounting for losses on a rotation rather than annual basis
\parencite{Zhao2016}, the paucity of research on desirable rotation crops in
two-crop systems \parencite{Pasley2021}, partial conversion of cropland to
pasture or wetland to meet nitrate targets \parencite{Randall2001,Gowda2007,Nangia2008},
and explicit joint evaluation of controlled drainage, rotation, cover crops, and
N management \parencite{Ma2007b}. Novel options---triticale with red clover
\parencite{Nance2007}, camelina--soybean relay-cropping \parencite{Emmett2022},
and fertilised, harvested rye for revenue \parencite{Malone2018,Malone2023}, and
diverse systems benchmarked against the two-year rotation \parencite{Gentry2025}
---are all advanced. \textcite{Dinnes2002} call broadly for perennial cash-cropping
systems that maintain returns while reducing leaching, and \textcite{Jaynes2001}
conclude that dramatic management change on artificially drained soils will be
required before corn production can be considered sustainable.